\section{格点研究概述}

我们在超立方格点上使用普通的单方块（plaquette）威尔逊作用量。我们利用扭转约化（twisted reduction）\cite{twist8}，它使得在大小为 $V=7^4$ 的小格点上可以实现一段相当不错的格距区间。为使扭转最佳地起作用，需要 $N$ 为 $4$ 乘一个素数。我们主要在 $N=44$ 下工作，并在 $N=28$ 检验有限 $N$ 效应。我们的模拟方法采用热浴更新与过弛豫更新相结合，平衡可在合理的机时内达到。在我们的大多数分析中可以忽略一切统计误差。与任何少有先例可循的格点工作一样，工作量的很大一部分在于确定进入数值计算的众多参数的有用范围，以便能够处理相关的物理问题。贯穿全文，我们用 $b$ 表示格点规范耦合常数，它是裸 't Hooft 耦合的倒数；相对于常规格点耦合 $\beta$，$b$ 已经把 $N$ 的正确幂次提取出来：
\be
b=\frac{\beta}{2N^2}=\frac{1}{g_{\rm YM}^2N}\ee
为使叙述清楚，我们忽略一个技术细节：对于扭转模拟，$b$ 取为未扭转情形下取值的负值；我们将引用常规的正的 $b$ 值。

\subsection{涂抹的格点威尔逊圈}

我们只考虑限制在超立方格点一个平面内的矩形圈，而且大部分工作针对方形圈。不过，连续统的处理会关心包含任意形状光滑圈的那一类可观测量。具体地说，连续统中的显式计算最好应用于位于 $R^4$ 中一个平坦平面内的圆形圈，因为它们具有最大对称性。


为了消除周长发散和拐角发散 \cite{div9}，我们构造与曲线相联系的威尔逊圈算符时，不使用由单方块威尔逊作用量的概率分布产生的 $SU(N)$ 链矩阵 $U_\mu (x)$，而使用它们的涂抹版本。我们采用 APE 涂抹 \cite{ape10}，从 $n=0$ 出发递归定义，此时涂抹后的矩阵等于 $U_\mu(x)$。令 $\Sigma_{U^{(n)}_\mu (x;f)}$ 表示与链 $U^{(n)}_\mu(x;f)$ 相联系的``staple''（三边组合），它由整套 $U^{(n)}_\nu(y;f)$ 矩阵表出。递归的一步把集合 $U^{(n)}_\mu (x;f)$ 变为集合 $U^{(n+1)}_\mu (x;f)$：
\begin{eqnarray}
X^{(n+1)}_\mu (x;f)= (1-|f|) U^{(n)}_\mu (x;f)+\frac{f}{6}
\Sigma_{U^{(n)}_\mu (x;f)}\nonumber\\ U^{(n+1)}_\mu (x; f
)=X^{(n+1)}_\mu (x;f) \frac{1}{\sqrt{[X^{(n+1)}_\mu (x;f)]^\dagger
X^{(n+1)}_\mu (x;f)}}
\end{eqnarray}
在模拟中从未出现上述方程中的幺正投影因 $X^{(n)}$ 奇异而停滞的情形。换言之，涂抹以概率一的方式是良定义的。对于扭转模拟，$f$ 取为未扭转情形取值的负值；我们将引用常规的正的 $f$ 值。

$U^{(n)}_\mu (x;f)$ 在规范变换下的变换方式与 $U_\mu (x)$ 相同。我们的涂抹威尔逊圈算符 $\hat W [L_1,l_2 ; f ; n]$ 定义为绕限制在一个平面内的 $L_1 \times L_2$ 矩形的有序乘积。$L_\alpha$ 是整数，给出以格距为单位的圈尺寸。当被走过的链从格点 $x=(x_1,x_2,x_3,x_4)$，$ x_\mu\in Z$ 出发并连接到 $\mu$ 正方向上的相邻格点 $x+\mu$ 时，链矩阵为 $U^{(n)}(x;f)$；而当该定向链沿反方向走过时，链矩阵为 $U^{\dagger (n)}(x;f)$。$\hat W$ 依赖于圈在哪里被剪开，但其本征值不依赖。本征值的集合在作用于 $U_\mu (x)$ 的基本规范变换下不变。

我们通过数值手段（后文将更详细描述）确认：当 $L_1 , L_2 , f, n$ 变化时，在特定的格点耦合下，$\hat W [L_1 , L_2 ,; f ; n] $ 的谱对非常小和/或涂抹很重的圈张开一个能隙。该能隙对非常大和/或涂抹很轻的圈闭合。


\subsection{格点大 $N$ 相变在连续统极限中存活}

即使只限于方圈，我们的格点算符也依赖许多参数：有一个整体标度，即方形的边长 $L$；有一个规范耦合 $b$，它固定物理单位下的格距、因而也控制标度；还有两个参数 $f$ 和 $n$ 控制 APE 涂抹，并同时负责驯服算符之迹在所有不可约 $SU(N)$ 表示中的周长发散和拐角发散。我们的构造隐式地调整了与每个表示中迹的重整化相联系的无穷多个未定有限部分，使得整个集合可以被看作来自一个（涨落的）$SU(N)$ 矩阵。这保证了圈的回折片段对算符没有贡献。

显然，我们在定义正则化威尔逊圈算符时所做选择中含有相当大的任意性。若把构造保持得更一般，这种任意性会在连续统中以一个大变换集合下不变性的形式反映出来，或许表现为一个广义重整化群方程。

我们的第一个目标是弄清如何调节 $\hat W$ 中的参数依赖，使得在格点上识别出的、能隙恰好张开处的 $N=\infty$ 相变点能在连续统极限中存活；在该极限中，格点耦合 $b$ 与 $L_{1,2}$ 一同趋于无穷，同时物理长度 $l_\alpha = L_\alpha a(b)$ 保持固定。$a(b)$ 是耦合 $b$ 处的格距，我们用物理的有限温退禁闭相变温度 $t_c$ 把它表示为
\be
a(b)=\frac{T_c(b)}{t_c}
\ee
其中 $\frac{1}{T_c(b)}$ 是 $S^1\times R^3$ 格点上紧致方向的格点尺寸，对应于有限温相变的临界耦合 $b$。

我们把涂抹步数 $n$ 取为与周长的平方成正比（我们把圈的尺寸限制为偶数 $L_1+L_2$），即 $n=\frac{(L_1+L_2)^2}{4}$。圈的物理尺寸为 $l_\alpha$ 且 $L_\alpha = \frac{l_\alpha}{a(b)}$。
我们取 $n=\frac{(L_1 (b) + L_2 (b))^2}{4}$，因为从物理上看乘积 $f n$ 是一个长度的平方。这一点可以从下一小节给出的简单微扰论证看出，或者更直观地，注意到涂抹是一个把圈加粗的无规行走、其厚度随涂抹步数的平方根增长。我们对 $n$ 的选取使 $f$ 在物理意义上成为无量纲参数；而在格点上 $f$ 实际上被限制在一个量级为一的区间内。

我们分析了 $\hat W [L_1 , L_2; f;
n=\frac{(L_1 + L_2 )^2}{4}]$ 的四个序列，每个序列保持物理尺寸固定在四个值之一。我们取 $L_1=L_2=L$。
通过各种手段（后文描述），我们改变 $f$ 并找到每个序列中每个成员恰在其谱中于本征值 $-1$ 处张开能隙的 $f$ 临界值。
（只要 CP 不变性得到保持，对任何圈，
本征值概率在 $\pm 1$ 处都是极值的。）然后把每个序列的临界值集合外推到 $b\to\infty$ 极限。在每个序列中，圈的物理尺寸按定义保持固定。这样，对我们选定的四个不同物理圈尺寸得到了四个连续统的 $f$ 临界值。自始至终，所用的 $f$ 值都在合理范围内，远在格点允许的最大范围 $0<f<0.75$ 之内 \cite{pertsmear11}。
事实上能够安排得使大 $N$ 相变附近只需要如此合理的 $f$ 值，这是我们的最基本结果，也是平面纯杨—米尔斯理论动力学的一个非平凡特征。以其物理尺度标识的每个序列相联系的相变点处 $f$ 临界值的连续统极限的全体，定义了一条临界线 $f_c(l)$。
对任何有限的 $b$ 都存在有限的格距效应，但它们足够小，使得 $f$ 的所有格点近似都保持在良好范围内。

于是我们得出结论：我们已在平面 QCD 中找到一条连续统临界线 $f_c(l)$，在其上我们的连续统威尔逊圈算符具有良定义的谱，而且该谱是临界的，意义是其本征值密度 $\hat \rho (\theta)$ 对任何 $\theta\ne\pi$ 非零但 $\hat\rho (\pi)=0$。这里我们把威尔逊圈算符的本征值记为 $e^{\imath\theta}$，其中 $\theta$ 是圆上的角度。

对于 $\Lambda l_\alpha >>1$，无必要保持 $n$ 与 $L_\alpha$ 二次相关；具有临界 $f$ 的圈是那些 $\Lambda l $ 为单位量级的尺寸 $l$。对大得多的圈，应当允许 $f n$ 变得与量级为 $\frac{1}{\Lambda^2}$ 的物理标度成正比，而不随物理圈尺寸趋于无穷而继续增长。要紧的是：对量级为 $\frac{1}{\Lambda}$ 的圈尺寸按上述方式调节 $n$，而对小圈让 $f n$ 也变小。对小圈我们希望进入微扰区域，因而不应在短距离上消除紫外涨落，因为在那里紫外涨落占主导。

一旦 $f_c (l)$ 线的存在在连续统极限中被确立，且 $\frac{df_c}{dl} > 0$，我们就可以比较有把握地确信：胀大一个 $\hat W [L,L;f;n]$ 使其在格点耦合 $b$ 处经过点 $L=\frac{l}{a(b)},
f=f_c (l)+{\cal O}(a^2(b)), n=L^2$，则该 $\hat W$ 的谱将经历相变。我们仍需确保：能隙一旦闭合，$-1$ 处的密度便从零开始连续建立起来；这将意味着相变是连续的，并很可能具有普适品格。如果相变是不连续的，例如 Polyakov 圈算符的有限温相变那样，那么相变两侧不太会存在标度区域，连接两侧的单一标度区域（我们所希望的那种类型）就被排除。有限温情形显然正是如此，因为那里的相变并非由取 $N\to\infty$ 引起，而是对所有 $N$ 都存在，且从 $N=3$ 起就是一级相变。


我们通过沿物理 $(f,l)$ 平面中另一条线开展附加模拟来解决这一点；该线不是常数 $l$，而是让 $l$ 在相当大的幅度上变化。我们选取一段方便的长度区间 $l$ 以及 $(f,l)$ 平面中的一条线 $f=f_0-f_1 l^2$，常数 $f_1 >0$ 的选取使得 $f$ 在 $f_1$ 与 $f_2$ 之间变化，其中 $f_1\approx
0.1$，$f_2\approx 0.2$。这条线以不太接近零的角度穿过 $f_c (l)$，从而当把 $\hat W$ 沿线在两个极端长度之间移动时增强了谱的可变性。这第二条线本身并无特别的物理含义，只是沿着它圈被胀大，同时涂抹量减少；两种效应都有助于闭合小圈时存在的谱隙。

为沿这条连续统线行进，我们固定格点圈尺寸 $L$ 而改变 $b$。我们同样保持 $n=L^2$ 固定，但按公式 $f=f(b)=f_0-f_1 (La(b))^2$ 改变 $f$。
沿这条线跟踪 $\hat W$，我们确认：当收敛到 $f_c(l)$ 的格点曲线带被穿过时，$\hat W$ 的谱确实张开一个能隙。

我们避免在靠近相变处工作，并用各种外推方法在两侧独立估计相应各相终止之处：在有能隙的一侧，我们研究能隙的缩小以估计它何时闭合。在无隙一侧，我们估计 $\hat \rho (\pi)$ 何时为零。

我们对三个递增的 $L$ 值重复了这一程序。
每种情形覆盖相似（重叠但不完全相同）的物理尺寸 $l$ 范围。在范围的公共部分，格点 $L$ 值越大，格点效应理应越小。我们发现总体上与按格距平方进行的标度破坏一致。

在每个相中分别进行的外推给出了各相端点的独立估计，它们彼此相当接近。这与 $\hat W$ 的本征值分布张开能隙的连续大 $N$ 相变相一致。从两侧分别确定的临界点的排序方式与不连续相变不相容；在不连续相变中 $\hat \rho(\pi)$ 应先降到某个非零值然后跳到零，跃入在 $\pi$ 附近具有有限能隙的相。于是我们断定：相变是连续的，并且可以指望在相变点附近有某种普适描述。实际上，很难想象何种动力学能够在谱中产生一级的大 $N$ 相变。

\subsection{涂抹消除紫外发散}

涂抹的效应在微扰论中容易理解：假定涂抹迭代的每一步都可以线性化。写出 $U^{(n)}_\mu
(x;f)=\exp(iA^{(n)}_\mu (x;f))$ 并对 $A_\mu$ 展开，人们在格点傅里叶空间中得到~\cite{pertsmear11}：
\be
A^{(n+1)}_\mu (q;f)= \sum_\nu h_{\mu\nu} (q) A^{(n)}_\nu (q;f) \label{eq:smearA}\ee
其中
\be
h_{\mu\nu} (q)= f(q)(\delta_{\mu\nu} - \frac {{\tilde q}_\mu
{\tilde q}_\nu}{\tilde q^2})+\frac {{\tilde q}_\mu {\tilde
q}_\nu}{\tilde q^2} \label{eq:hkern}\ee 这里 $\tilde q_\mu = 2\sin
(\frac{q_\mu}{2})$，且
\be
f(q)=1-\frac{f}{6}\tilde q^2
\label{eq:fq}
\ee
迭代的解法是把 $f(q)$ 换成 $f^n (q)$，对足够小的 $f$，
\be
f^n(q) \sim e^{-\frac{f n}{6} \tilde q^2}
\label{eq:fnsol}
\ee

条件 $0<f(q)<1$（对所有动量）要求 $0<f<0.75$。为对数字有些感觉：一个 $4\times 4$ 圈，在使其临界的 $b$ 值处将有 $n=16$ 和 $f\sim 0.15$，这使得 $\frac{f n}{6}\sim 0.4$，即在长达圈子本身的距离上的加粗量。这就是为什么大得多的圈不应该用以周长平方方式持续增长的 $f n$ 因子来涂抹；毋宁说，例如对边长为 $L$ 的方圈，如下选取将是适当的：
\be
f=\frac{\tilde f}{1+M^2 L^2}
\ee
这里 $M$ 是以格点单位表达的物理质量 $m$，$M=m a(b)$。$m$ 是典型的强子尺度，选取它使得在我们发现的大 $N$ 相变处 $ML$ 小于比如说 $0.01$。

为了让我们确信 $f$ 的这种截断依赖扩展确实能与显示禁闭、且弦张力取已知值的大圈光滑衔接，我们在所研究范围的大尺寸一端用涂抹圈粗略估计了弦张力。我们发现得到的弦张力接近裸格点威尔逊圈的值，这表明我们的涂抹威尔逊圈算符是可接受的纯 $SU(N)$ 规范理论弦候选者。


总之，涂抹对 $A^{(n)}_\mu (x;f)$ 传播子的效应是消除引起裸威尔逊圈算符迹中周长发散与拐角发散的紫外奇异性。当逼近连续统极限时乘积 $fn$ 以物理长度平方的方式变化，这就奏效。对小圈我们选这个物理长度正比于圈子本身。对大圈，物理长度稳定在某个强子尺度。这样一来，我们的算符在标度非常小时对短距离动力学敏感；而当禁闭成为主导效应、圈非常大时，涂抹威尔逊圈的弦张力就是该理论真实、普适的弦张力，不受涂抹影响。当圈非常大时，描述威尔逊圈形状依赖的有效弦理论的领头项预期是普适的，弦张力仅以量纲单位的平凡方式进入。

我们要强调，我们并不是在执行重整化群变换；我们的威尔逊圈算符定义为时空中的无限细圈、并对所有圈尺寸成立；毋宁说，我们在应用一种正则化，其目标是使威尔逊圈算符在连续统中有限。如果我们对每个不可约表示单独处理，消除周长发散和拐角发散将留下无穷多个任意的有限部分。通过施加威尔逊圈算符幺正性的要求，这一自由度被隐式利用，它保证了 Polyakov 所谓的``之字形''（zigzag）对称性 \cite{poly3}。保持威尔逊圈算符的幺正品格又反过来使本征值分布的概念在连续统中良定义，并为可能出现的大 $N$ 相变搭好了舞台。之字形对称还影响无穷大 $N$ 处格点圈方程的形式，基础表示中裸威尔逊圈的迹服从这些方程；大圈的有效弦理论不必服从之字形对称，但这种对称可以对其耦合常数提供有用的约束。
